\documentclass[../../main.tex]{subfiles}
\begin{document}

\begin{flushright}
\footnotesize\textit{【自動文字起こし・要確認】}
\end{flushright}

{\bf [解]}
$a,b>0$とし，$f(x)=x^a$，$g(x)=\log bx$（$x>0$）とおく．$f'(x)=ax^{a-1}$，$g'(x)=\dfrac{1}{x}$である．接点を$P(s,t)$とすると，$f(s)=g(s)$，$f'(s)=g'(s)$がともに成り立つ：
\begin{align}
s^a=\log bs=t,\qquad as^{a-1}=\dfrac{1}{s}.
\end{align}
第2式の両辺に$s$をかけて$as^a=1$，すなわち$s^a=\dfrac{1}{a}$（$s>0$）だから$s=\bigl(\dfrac{1}{a}\bigr)^{1/a}$．第1式に代入すると
\begin{align}
t=s^a=\dfrac{1}{a},\qquad b=\dfrac{1}{s}e^{1/a}=(ae)^{1/a}.
\end{align}

\textbf{(2)} $y=g(x)$，$y=f(x)$のグラフは，$0<x<s$で$f(x)>g(x)$（$g(x)\to-\infty\ (x\to0^+)$に対し$f(x)\to0$）．$0<h<\dfrac{1}{a}$のもとで考えれば十分で，
\begin{align}
A(h)=\int_h^s\bigl(f(x)-g(x)\bigr)dx=\Bigl[\dfrac{1}{a+1}x^{a+1}-x(\log b+\log x-1)\Bigr]_h^s.
\end{align}
これを展開すると
\begin{align}
A(h)=\dfrac{1}{a+1}\bigl(s^{a+1}-h^{a+1}\bigr)-s(\log b+\log s-1)+h(\log b+\log h-1).
\end{align}
$h\to0^+$のとき$h^{a+1}\to0$（$a+1>0$），$h\log h\to0$より$h(\log b+\log h-1)\to0$なので
\begin{align}
\lim_{h\to0}A(h)=\dfrac{1}{a+1}s^{a+1}-s(\log b+\log s-1).
\end{align}
ここで$\log b+\log s=\dfrac{1}{a}(\log a+1)+\Bigl(-\dfrac{1}{a}\log a\Bigr)=\dfrac{1}{a}$だから$\log b+\log s-1=\dfrac{1-a}{a}$．また$s^{a+1}=s^a\cdot s=\dfrac{s}{a}$．よって
\begin{align}
\lim_{h\to0}A(h)=\dfrac{s}{a(a+1)}-s\cdot\dfrac{1-a}{a}=\dfrac{s}{a}\Bigl(\dfrac{1}{a+1}-(1-a)\Bigr)=\dfrac{s}{a}\cdot\dfrac{1-(1-a)(a+1)}{a+1}=\dfrac{s}{a}\cdot\dfrac{a^2}{a+1}=\dfrac{a}{a+1}s.
\end{align}
$s=\bigl(\dfrac{1}{a}\bigr)^{1/a}$だから
\begin{align}
\lim_{h\to0}A(h)=\dfrac{a}{a+1}\Bigl(\dfrac{1}{a}\Bigr)^{1/a}.
\end{align}


\newpage

\end{document}
